\documentclass{beamer}
\usepackage{ctex}
\usetheme{SimplePlus}
\setbeamertemplate{footline}{}

\usepackage{amsmath,amssymb,mathtools}
\usepackage{tikz}
\usetikzlibrary{arrows.meta}
\usepackage{tikz-cd}

\usepackage{xcolor}

\definecolor{myblue}{RGB}{40,40,150}

\usepackage{unicode-math}
%\setmathfont{STIX Two Math}
%\setmathfont{TeX Gyre Termes Math}
%\setmathfont{Libertinus Math}
%\setmathfont{Fira Math}

\usepackage[
  hybrid,                            % 允许在 Markdown 中混用 LaTeX 命令
  texMathDollars                     % 启用 $...$ 和 $$...$$ 数学公式
]{markdown}

\usepackage{graphicx} % Allows including images
\graphicspath{{./images/}}

\title{comprehension categories 相关概念}

\date{}

\begin{document}

\begin{frame}
    % Print the title page as the first slide
    \titlepage
\end{frame}

\begin{frame}[fragile]{cartesian morphism\cite{ncatlab_grothendieck_fibration}}
Given a functor $P \colon \mathcal{E} \to \mathcal{B}$, a morphism $f \colon x \to y$ in its domain category $\mathcal{E}$ is called \textit{cartesian} if for any morphism $g \colon z \to y$ in $\mathcal{E}$ and $w \colon P(z) \to P(x)$ in $\mathcal{B}$ such that $P(f) \circ w = P(g)$, there exists a \textit{unique} $\widehat{w} \colon z \to x$ such that $g = f \circ \widehat{w}$ and $P(\widehat{w}) = w$:

$$
\begin{tikzcd}[row sep=large, column sep=large]
 & z \arrow[dr, dashed, "\exists! \, \widehat{w}"'] \arrow[drr, bend left=15, "\forall g"] & & \\
\mathcal{E} \arrow[dd, "P"'] & & x \arrow[r, "f", "\text{cartesian}"'] & y \\
 & P(z) \arrow[dr, "\forall w"'] \arrow[drr, bend left=15, "P(g)"] & & \\
\mathcal{B} & & P(x) \arrow[r, "P(f)"] & P(y)
\end{tikzcd}
$$
\end{frame}

\begin{frame}[fragile]{Cartesian morphism 的直觉}
\begin{markdown}
- $f$ 是沿底空间 $P(f)$ 在总空间中的 **“主拉回”**
- cartesian morphism 是 category-theoretic version of substitution/reindexing。
\end{markdown}
\end{frame}

\begin{frame}{Grothendieck fibration}
A functor $P \colon \mathcal{E} \to \mathcal{B}$ is a \textit{fibration} if for all objects $y \in \mathcal{E}$ and morphisms $f_0 \colon x_0 \to P(y)$, there is a cartesian morphism $f \colon x \to y$ such that $P(f) = f_0$.

Such $f$ is also called a ``cartesian lifting'' of $P(f)$ to $y$, and a choice of cartesian lifting for every $y$ and $P(f)$ is called a \textit{cleavage}.
\end{frame}

\begin{frame}[fragile,shrink]{Grothendieck 纤维化}
\begin{markdown}
函子
$$
p:\mathcal E\to\mathcal B
$$

称为一个 **Grothendieck fibration**，如果满足：

对于任意对象 $Y\in\mathcal E$ 和任意基底态射
$$
u:I\to pY
$$

都存在一个笛卡尔态射
$$
\overline u:X\to Y
$$

使得
$$
p(\overline u)=u.
$$

也就是说，每个对象 $Y$ 都可以沿着任意基底态射 $u:I\to pY$ 被“反向运输”到纤维 $\mathcal E_I$。

常把 $\overline u$ 写成：
$$
u^*Y\longrightarrow Y.
$$

其中 $u^*Y\in\mathcal E_I$。

这个 $u^*Y$ 可以理解为 $Y$ 沿 $u$ 的拉回、重索引或替换。
\end{markdown}
\end{frame}

\begin{frame}[fragile]{笛卡尔提升不要求唯一选择}
\begin{markdown}
Grothendieck fibration 只要求：

> 这样的笛卡尔提升存在。

但它不要求我们为每个 $(u,Y)$ 指定一个具体的提升。

事实上，两个不同的笛卡尔提升之间存在唯一的垂直同构。也就是说，如果
$$
X\to Y
\quad\text{和}\quad
X'\to Y
$$

都是 $u:I\to pY$ 的笛卡尔提升，那么存在唯一的同构
$$
X\cong X'
$$

位于纤维 $\mathcal E_I$ 中，并且与到 $Y$ 的两个态射相容。

因此，反向拉回 $u^*Y$ 通常只在“唯一同构”的意义下确定，而不是严格唯一。

这正是为什么一般得到的是伪函子，而不是严格函子。
\end{markdown}
\end{frame}

\begin{frame}[fragile]{Cleavage：劈分、笛卡尔提升的选择}
\begin{markdown}
设 $p:\mathcal E\to\mathcal B$ 是一个 fibration。

一个 **cleavage** 是对每一个：

- 对象 $Y\in\mathcal E$；
- 基底态射 $u:I\to pY$；

选择一个笛卡尔态射
$$
\overline u_Y:u^*Y\to Y
$$

满足
$$
p(\overline u_Y)=u.
$$
换句话说，cleavage 就是：

> 为所有需要的笛卡尔提升，选定一个代表。

通常要求恒等态射对应于恒等笛卡尔态射：
$$
(1_{pY})^*Y=Y,
\qquad
\overline{1_{pY}}_Y=1_Y.
$$
但不同文献对 cleavage 是否自动包含单位条件，约定略有不同。

\end{markdown}
\end{frame}

\begin{frame}[fragile]{为什么需要 cleavage？}
\begin{markdown}
没有 cleavage 时，对每个 $u:I\to J$ 和 $Y\in\mathcal E_J$，只能说存在一个对象 $u^*Y$，并且它在同构意义下唯一。

有了 cleavage 后，才可以写出一个具体的操作：
$$
u^*:\mathcal E_J\to\mathcal E_I.
$$

对纤维内态射 $f:Y\to Y'$，笛卡尔泛性质给出一个唯一的态射
$$
u^*f:u^*Y\to u^*Y'.
$$

于是得到反变的“重索引函子”：
$$
u^*:\mathcal E_J\to\mathcal E_I.
$$
\end{markdown}
\end{frame}

\begin{frame}[fragile]{Cloven fibration：带劈分的纤维化}
\begin{markdown}
一个 **cloven fibration** 就是：
$$
p:\mathcal E\to\mathcal B
$$

再加上一个 cleavage。

所以：
$$
\boxed{\text{cloven fibration} = \text{fibration} + \text{一个指定的 cleavage}}
$$

必须注意：

- fibration 只要求笛卡尔提升存在；
- cloven fibration 还指定了具体选择；
- 同一个 fibration 可以有很多不同的 cleavage；
- cleavage 通常不唯一，也不一定严格地与复合相容。
\end{markdown}
\end{frame}

\begin{frame}[fragile]{Split fibration：严格可复合的劈分}
\begin{markdown}
一个 cleavage 称为 **split cleavage**，如果满足严格的：
$$
(1_I)^*=\operatorname{id}_{\mathcal E_I}
$$

以及对于
$$
I\xrightarrow{u}J\xrightarrow{v}K,
$$

有
$$
(vu)^*=u^*v^*
$$

严格成立，而不只是自然同构。

因此：
$$
\boxed{\text{split fibration} = \text{cloven fibration} + \text{严格满足单位和复合}}
$$

一般 cleavage 只给出自然同构：
$$
u^*v^*Y\cong (vu)^*Y.
$$

所以一般 cloven fibration 对应的是伪函子，而 split fibration 对应严格函子。
\end{markdown}
\end{frame}

\begin{frame}[fragile]{Comprehension Categories\cite{jacobs_comprehension}}
\begin{enumerate}
  \item a category $\mathcal{C}$,
  \item a (cloven) fibration $p \colon \mathcal{T} \to \mathcal{C}$,
  \item a functor $\chi \colon \mathcal{T} \to \mathcal{C}^\to$ preserving cartesian arrows,
\end{enumerate}
\noindent such that the following diagram commutes.

\[
\begin{tikzcd}[column sep=large, row sep=large]
\mathcal{T} \arrow[rr, "\chi"] \arrow[dr, "p"'] & & \mathcal{C}^\to \arrow[dl, "\text{cod}"] \\
& \mathcal{C} & 
\end{tikzcd}
\]

\noindent A comprehension category is
\begin{itemize}
  \item[\textbf{full}] if $\chi$ is full and faithful;
  \item[\textbf{split}] if $p$ is a split fibration.
\end{itemize}
\end{frame}

\begin{frame}[fragile]{Interpreting MLTT in a Comprehension Category}
\[
\begin{tikzcd}[column sep=large, row sep=large]
\mathcal{T} \arrow[rr, "\chi"] \arrow[dr, "p"'] & & \mathcal{C}^\to \arrow[dl, "\text{cod}"] \\
& \mathcal{C} & 
\end{tikzcd}
\]

\vspace{1em}

\begin{tabular}{rl}
\textcolor{myblue}{Category $\mathcal{C}$} & models contexts \\
\textcolor{myblue}{Fibre $\mathcal{T}_\Gamma$} & models types in context $\Gamma$ \\
\textcolor{myblue}{Reindexing} & models substitution \\
\textcolor{myblue}{Comprehension $\chi$} & models context extension $(\Gamma, A) \mapsto (\Gamma.A \xrightarrow{\chi(A)} \Gamma)$ \\
\textcolor{myblue}{Sections of $\chi(A)$} & model terms $\Gamma \vdash t : A$ \\
\end{tabular}
\end{frame}


% 参考文献帧，长列表自动分页
\begin{frame}[allowframebreaks]{参考文献}
\begin{thebibliography}{99}

\bibitem{ncatlab_grothendieck_fibration}
\url{https://ncatlab.org/nlab/show/Grothendieck+fibration}

\bibitem{jacobs_comprehension}
Bart Jacobs. “Comprehension Categories and the Semantics of Type Dependency”. In: Theor. Comput. Sci. 107.2 (1993), pp. 169–207.

\end{thebibliography}
\end{frame}
\end{document}